% tikz-tensors-edges.tex -- edge types, and the routes an edge takes.
%
% An edge type is a TikZ style built on `tn edge': what the line looks like.
%   tn edge    plain (with a colour: tn edge=exact)
%   tn wavy    a wave, undecorated where it has to turn
%   tn arrow   an arrowhead part way along; `tn arrows=false' turns every one
%              off, and which way it points is the route's to say
%   tn double  two indices fused into one
% A user's own type is a style of the same kind. Where an edge goes is not the
% type's business: a route below takes any type, and the layouts choose the
% route.
\tikzset{
  tn wavy/.style   = {tn edge, decorate,
                      decoration={snake, amplitude=0.55mm,
                                  segment length=2.4mm, pre length=0.4mm,
                                  post length=0.8mm}},
  tn arrow/.style  = {tn edge, postaction={decorate},
                      decoration={markings, mark=at position 0.58 with
                        {\iftnarrows\arrow{Latex[length=1.7mm]}\fi}}},
  tn double/.style = {tn edge, double=grey1, double distance=1.2pt},
}

% ---- routes -----------------------------------------------------------------
% Every route is drawn by \tnbond, so under the nodes. Coordinates are lengths
% (the layouts work in them), x before y.
%
% \tn@line{<style>}{<x1>}{<y1>}{<x2>}{<y2>}     straight, from 1 to 2
\def\tn@line#1#2#3#4#5{\tnbond[#1]{(#2,#3) -- (#4,#5)}}
% \tn@turns{<style>}{<path>}   a path of straight runs, its corners rounded as
% little as reads as one line turning rather than two lines meeting
\def\tn@turns#1#2{\tnbond[#1, rounded corners=1.5pt]{#2}}
% \tn@bond{<style>}{<x left>}{<x right>}{<y>}{<points of left>}{<of right>}
% Straight along a layer, its direction taken from the shapes it joins: away
% from one pointing right, into one pointing left, so that in a canonical form
% every arrow runs toward the centre. Anything else is drawn left to right.
\def\tn@bond#1#2#3#4#5#6{%
  \edef\tn@pa{#5}\edef\tn@pb{#6}%
  \def\tn@rev{0}%
  \ifx\tn@pb\tn@sleft \ifx\tn@pa\tn@sright\else\def\tn@rev{1}\fi\fi
  \ifx\tn@pa\tn@sleft \ifx\tn@pb\tn@sright\else\def\tn@rev{1}\fi\fi
  \ifnum\tn@rev=1 \tn@line{#1}{#3}{#4}{#2}{#4}%
  \else \tn@line{#1}{#2}{#4}{#3}{#4}\fi}

% \tnswap[<style>]{<left top>}{<right top>}{<drop>}
% Two legs that start at <left top> and <right top> and end <drop> lower,
% exchanged: short straight stubs in <style> at both ends, joined by a smooth
% crossing in the same style undecorated (a wavy line on a diagonal reads as
% noise).
\newcommand{\tnswap}[4][tn edge]{%
  \begin{pgfonlayer}{tnedges}%
    \path #2 coordinate (tn@l) #3 coordinate (tn@r);
    \draw[#1] (tn@l) -- ++(0,-0.3) coordinate (tn@l1);
    \draw[#1] (tn@r) -- ++(0,-0.3) coordinate (tn@r1);
    \draw[#1] ($(tn@l)+(0,-#4)$) -- ++(0,0.3) coordinate (tn@l2);
    \draw[#1] ($(tn@r)+(0,-#4)$) -- ++(0,0.3) coordinate (tn@r2);
    \draw[#1, decorate=false] (tn@l1) .. controls +(0,-0.3) and +(0,0.3) .. (tn@r2);
    \draw[#1, decorate=false] (tn@r1) .. controls +(0,-0.3) and +(0,0.3) .. (tn@l2);
  \end{pgfonlayer}%
}
